\subsection{Tổng quan mô hình tin cậy}

Hệ thống bỏ phiếu điện tử được thiết kế theo nguyên tắc \textit{defense in depth} — không tin cậy bất kỳ tầng đơn lẻ nào. Mỗi tầng bảo vệ vận hành độc lập, sử dụng cơ chế kỹ thuật khác nhau, do đó việc một tầng bị xâm phạm không đủ để phá vỡ tính toàn vẹn của cuộc bầu cử. Toàn bộ kiến trúc bảo mật được tổ chức thành bảy tầng phòng vệ theo thứ tự từ ngoài vào trong, kết hợp giữa các kỹ thuật bảo mật ở tầng giao thức, tầng ứng dụng và tầng dữ liệu:

\begin{itemize}
    \item \textbf{Tầng 1 — Bảo mật HTTP:} bao gồm Helmet (đặt các tiêu đề bảo mật chuẩn), CORS (giới hạn nguồn gốc gọi API), và Rate Limiting (chống tấn công làm tràn yêu cầu).

    \item \textbf{Tầng 2 — Xác thực và phân quyền JWT:} dùng cặp khóa bí mật riêng cho access token và refresh token, kèm cơ chế kiểm soát truy cập theo vai trò (RBAC).

    \item \textbf{Tầng 3 — Quản lý phiên trên Redis:} mỗi phiên bỏ phiếu được lưu với thời gian sống ngắn, gắn cờ \texttt{signed}/\texttt{voted} và được dùng làm máy trạng thái phía máy chủ.

    \item \textbf{Tầng 4 — Chữ ký mù tập thể EC-Schnorr:} bảo đảm bí mật phiếu bầu ngay cả khi máy chủ ký không trung thực và ngăn cử tri ký quá một lần.

    \item \textbf{Tầng 5 — Ràng buộc MongoDB:} chỉ mục duy nhất trên cặp \texttt{(electionId, voterId)} và \texttt{(electionId, blindedCommitment)}, cùng các giao dịch nhiều thao tác (multi-document transaction), ngăn lỗi đua tranh và phiếu trùng.

    \item \textbf{Tầng 6 — Cam kết Merkle:} cây Merkle gom toàn bộ phiếu vào một giá trị gốc 32 byte, được ghi bất biến lên sổ cái Fabric khi đóng bầu cử.

    \item \textbf{Tầng 7 — Sổ cái Hyperledger Fabric:} mạng riêng được cấp phép, có chứng chỉ X.509 cho mỗi nút, hỗ trợ chính sách xác nhận giao dịch (endorsement policy) và tính bất biến đảm bảo bởi liên kết băm khối liên tiếp.
\end{itemize}

Sáu tính chất an toàn của một hệ thống bầu cử điện tử được đảm bảo đồng thời nhờ sự kết hợp giữa các tầng nêu trên, được tổng hợp trong Bảng~\ref{tab:tinh-chat-bau-cu}.

\begin{table}[H]
\centering
\caption{Sáu tính chất bầu cử và cơ chế đảm bảo trong hệ thống}
\label{tab:tinh-chat-bau-cu}
\begin{tabularx}{\textwidth}{|L{0.4}|L{0.7}|}
\hline
\textbf{Tính chất} & \textbf{Cơ chế đảm bảo} \\
\hline
Tính bí mật (\textit{ballot secrecy}) & Chữ ký mù: máy chủ không biết cử tri bỏ phiếu cho ứng viên nào \\
\hline
Tính duy nhất (\textit{one voter -- one vote}) & Phiên Redis kết hợp chỉ mục duy nhất trên MongoDB và bộ khử trùng tại nút ký \\
\hline
Tính xác thực (\textit{vote authenticity}) & Kiểm tra chữ ký EC-Schnorr trong pha giải mù phiếu bầu \\
\hline
Tính không thể chối bỏ (\textit{non-repudiation}) & Sổ cái Fabric ghi bất biến mã giao dịch của mỗi phiếu \\
\hline
Tính kiểm chứng (\textit{verifiability}) & Bằng chứng Merkle (\textit{Merkle proof}) cho phép bất kỳ bên nào tự kiểm tra phiếu \\
\hline
Tính ẩn danh (\textit{unlinkability}) & Bảng \texttt{Vote} và \texttt{RevealedVote} không thể nối kết qua bất kỳ trường khóa nào \\
\hline
\end{tabularx}
\end{table}

Các tiểu mục tiếp theo của chương sẽ trình bày chi tiết từng tầng — bắt đầu từ lớp mật mã cốt lõi (EC-Schnorr), tiếp đến các tầng xác thực, phiên, blockchain và kết thúc bằng các kỹ thuật tối ưu hiệu năng giúp hệ thống đạt thông lượng cần thiết cho quy mô bầu cử thực tế.
